% ASSUMPTION (full Overleaf source unavailable): Section 2.2 = FEM data and their preparation; Section 2.3 = surrogate models and training protocol. Everything outside 3.2 is otherwise unchanged; this text replaces the former model/protocol subsections, including the "Adam followed by LBFGS ... cosine schedule" passage (ERRATA E-05: no such protocol exists in the code).
% Every number below carries a "% src:" comment naming the file it was read from. No v1 (pre-cleaning) figure is used.
% VERIFICATION PASS (2026-09-07): every number re-checked against publication/*.md, publication/code/*.py, publication/experiments/run_matrix.py and results/runs_v2.csv (341 rows, no error field). Corrections are marked "% FIXED:"; unverifiable items "% TODO:".
% IMPORTANT: the --hp JSON that produced n_layers=3, n_units=96, learning_rate=0.003 (and lambda_bc=0 for mlp) is NOT in the repository; runs_v2.csv logs only n_layers, n_units, learning_rate, lambda_physics, lambda_bc, phys_gain. Every other Config field quoted below (activation, weight decay, clipping, batch size, val_frac, n_fourier, n_test_funcs, n_quadrature, r_clip) is the code DEFAULT and cannot be confirmed from the log — see the TODO in "Training protocol".
% Cross-references: Fig.~\ref{fig:architecture} = publication/figures/fig_architecture.png (both tracks + loss table); Table~\ref{tab:splits} (Section 2.2, sec22_fem_dataset.tex) is the design whose results Table~\ref{tab:main} (tables/all_tables.tex) reports. % EDIT: the duplicate "Evaluation splits" table formerly in this file is deleted
\subsection{Surrogate models and training protocol}
\label{sec:models}

% REV К13: вводная конструкция перед описанием моделей ("При подготовке моделей машинного обучения для последующего сравнительного анализа использовалась ...")
% REV К13/К12 (правка ревизии): "a single set of training conditions ... for every family" было бы неверно — предел эпох, размер батча, patience и число сидов различаются между основным и расширенным треком (см. "Training protocol": 800 (300), 64 (32), 40 (30), сиды 0–4 / 0–2). Условия равны внутри трека; в перечень разделов добавлен последний абзац (К14).
In preparing the machine-learning models whose comparative analysis is reported
below, a single network architecture and a single optimiser were adopted for
every family, and within each track all families were trained under the same
conditions. The network and its inputs
are described first, then the two tracks and the treatment of the axial
coordinate, the loss function and its weak-form variant, the training protocol,
the metrics and the physical audit, the two backends, the statistical
treatment of the results and, last, the architecture and hyperparameters held
fixed throughout.

\paragraph{Network and inputs.}
% REV К12: термин "training budget" раскрыт (предел числа эпох, критерий ранней остановки с patience, размер мини-батча, сиды) и больше не стоит голым
% REV К12 (правка ревизии): снято утверждение "all ... under identical conditions ... the same set of random seeds" — оно противоречит абзацу "Training protocol" того же раздела (800 (300) эпох, батчи 64 (32), patience 40 (30), пять сидов 0–4 против трёх 0–2) и абзацу "Backends" (JAX-двойник берёт свой поток случайных чисел). Числа не добавлены и не изменены.
All surrogates share one network and one optimiser, and the families compared
with one another are trained under the same conditions: the same limit on the
number of epochs, the same early-stopping criterion with the same patience, the
same mini-batch size and batch-ordering rule, and the same seeds; the values,
which differ between the main and the extended track, are stated in full in
``Training protocol'' below. The
families differ only in the composition of the loss (Table~\ref{tab:families},
Fig.~\ref{fig:architecture}). In the main track the input is
$\mathbf{x}=(Q,k,\alpha,\mu,v,r)$, with $r\in[0,1]$ the radius normalised by the
free-surface radius of the same case ($r=1$ is the free surface).
% src: trainer.py (Net n_in=6, r last column); dataset.py (r = r_phys / r_phys[:, -1], r[-1] = 1)
The radius is augmented by a Fourier embedding $[\sin(2\pi k r),\cos(2\pi k r)]$,
$k=1,\dots,4$ \cite{TODO-fourier-features};
% src: trainer.py FourierR (k = 1..n, times 2π), Config.n_fourier=4 (default; not logged — see TODO below)
inputs and targets are standardised with training-set statistics; the body is a
multilayer perceptron with three hidden layers of 96 $\tanh$ units and a linear
head with four outputs
$\hat{\boldsymbol\sigma}=(\hat\sigma_{rr},\hat\sigma_{\theta\theta},\hat\sigma_{zz},\hat\tau_{rz})$,
de-standardised before entering any physical term.
% src: runs_v2.csv (n_layers=3, n_units=96 in all 341 runs); trainer.py Scalers.fit on the training subset; COMPONENTS, n_out=4
% TODO: tanh is the Config.activation default (ACTS = tanh/selu/softplus); activation is not logged in the csv and the hp JSON is absent — confirm with the author of the run.
All arithmetic is in double precision in both backends.
% src: trainer.py (self.dtype = torch.float64); jax_twin.py (jax.config.update("jax_enable_x64", True))
Each FEM case contributes 20 radial points, always kept in the same partition
(group-aware splitting).
% src: trainer.py N_R=20; splits.py sets_to_rows (set-major rows)

\paragraph{Two tracks and the role of $z$.}
% REV К15: абзац переписан по формулировке руководителя — уравнения равновесия вводятся как основа PINN, отмечена зависимость входящих в них величин от r и z, и указано, что при подготовке моделей соотношения (1)–(2) претерпевают ряд изменений
% TODO К15: подставить конкретные работы вместо \cite{TODO-pinn-source-1,TODO-pinn-source-2} — руководитель имеет в виду источники [17,18] прежнего списка литературы (постановка PINN для уравнений равновесия сплошной среды); в publication/ этих ссылок нет, нужен библиографический список от руководителя
The physics-informed families are built on the equations of equilibrium of an
axisymmetric body, which are given below. The quantities entering them depend
on both coordinates, $r$ and $z$. In preparing the models, however, following
\cite{TODO-pinn-source-1,TODO-pinn-source-2}, relations
\eqref{eq:eq_full_r}--\eqref{eq:eq_full_z} undergo a number of changes, which
are set out in the remainder of this paragraph.
In the main track the axial coordinate is not an input: each case is one radial
profile, the median along the steady-state axial window (25--75\,\% of the wire
length, $L/R\approx3.5$).
% src: QA_DATA_AND_Z.md §4 (axial window 25–75 %, median window length 57.5 mm at R = 16.2 mm, L/R ≈ 3.5); dataset.py (20 equally populated r-bins, median)
Axial derivatives of the network therefore do not exist, and the axisymmetric
equilibrium equations
\begin{align}
  R_r &= \frac{\partial\sigma_{rr}}{\partial r}
        + \frac{\partial\tau_{rz}}{\partial z}
        + \frac{\sigma_{rr}-\sigma_{\theta\theta}}{r} = 0, \label{eq:eq_full_r}\\
  R_z &= \frac{\partial\tau_{rz}}{\partial r}
        + \frac{\partial\sigma_{zz}}{\partial z}
        + \frac{\tau_{rz}}{r} = 0 \label{eq:eq_full_z}
\end{align}
% src: trainer2d.py docstring; dataset2d.py docstring; ERRATA.md E-02
can only be imposed with the axial terms dropped. Whether that is admissible was
measured on the raw FEM fields (2443 cases, $20\times20$ $(z,r)$ grid, median per
case): in \eqref{eq:eq_full_r} the dropped $|\partial\tau_{rz}/\partial z|$ is
1.3\,\% of the retained terms (75th percentile 3.9\,\%); in \eqref{eq:eq_full_z}
the dropped $|\partial\sigma_{zz}/\partial z|$ is 68\,\% (75th percentile 86\,\%).
% src: ERRATA.md E-02 table (0.013 / 0.039; 0.68 / 0.86; 2443 sets, grid 20 × 20, medians); QA_DATA_AND_Z.md §4 (1.3 % / 68 %)
The radial equation is therefore used in the reduced form
\begin{equation}
  R_1(r) = \frac{\partial\sigma_{rr}}{\partial r}
         + \frac{\sigma_{rr}-\sigma_{\theta\theta}}{r},
  \label{eq:eq_red}
\end{equation}
% src: trainer.py _radial_residual
whereas the axial equation is not used in the main track, for a reason stronger
than the size of the dropped term: the remainder is a total derivative,
\begin{equation}
  \frac{\partial\tau_{rz}}{\partial r} + \frac{\tau_{rz}}{r}
  \equiv \frac{1}{r}\,\frac{\partial\,(r\,\tau_{rz})}{\partial r},
  \label{eq:identity}
\end{equation}
so forcing it to zero means $r\,\tau_{rz}=\mathrm{const}$ and, with regularity on
the axis ($\tau_{rz}(0)=0$), $\tau_{rz}\equiv0$ over the whole radius. That is not
an approximate axial balance but the statement that shear is absent, which the
FEM data contradict (median relative range of $r\,\tau_{rz}$ along the radius
3.15).
% src: ERRATA.md E-02 (3.15); physics_check.py EQUILIBRIUM_R2_STATUS / zero_shear_residual docstring (3.15); trainer.py _zero_shear docstring (3.15)
The term exists in the code under its true name (zero-shear penalty,
$\lambda_{\mathrm{shear}}$) and is switched off, $\lambda_{\mathrm{shear}}=0$, in
every run reported here.
% src: trainer.py Config.lambda_shear=0.0, _zero_shear (added to the loss only if lambda_shear > 0)
% TODO: lambda_shear is not logged in runs_v2.csv; "every run" rests on the default and the absent hp JSON.

In the extended track the axial coordinate is added,
$\mathbf{x}=(Q,k,\alpha,\mu,v,z,r)$; the field is sampled on an $8\times20$ $(z,r)$
grid inside the same window, one median per cell, without collapsing the axial
direction; $z$ is normalised to the window (0 = inlet, 1 = outlet) and $r$ to the
surface radius of the same axial section.
% FIXED: "without axial averaging" → "one median per cell, without collapsing the axial direction" (dataset2d.grid_one takes the median within each cell)
% src: trainer2d.py (Net2D n_in=7, column order [Q,k,α,μ,v,z,r]); dataset2d.py (build n_z=8, n_r=20; z = (ZP − ZP[0])/span; r = RP / RP[:, :, -1]); QA_DATA_AND_Z.md §4; RESULTS.md §5 (8 × 20)
Both equations are then written in full, and three families are compared:
\texttt{mlp2d} (no physics), \texttt{pinn2d\_reduced} ($R_r$ without
$\partial\tau_{rz}/\partial z$, exactly the operator of the main track) and
\texttt{pinn2d\_full} (both residuals).
% src: trainer2d.py FAMILIES_2D, _loss_with_physics (full: d_tau_dz, d_szz_dz; reduced: zeros)
Physical derivatives follow from the chain rule,
$\partial/\partial r_{\mathrm{phys}}=R^{-1}\partial/\partial r$,
$\partial/\partial z_{\mathrm{phys}}=L^{-1}\partial/\partial z$, with $R$ and $L$
the surface radius and window length of the case; the $1/r$ factor uses $r$
clipped at $0.05R$, since nodal radii near the axis carry noise (worst
$|r|/R=0.054$).
% src: trainer.py (_physics_strong: dsrr_dnorm / R; r_safe = max(r_phys, R·r_clip)); trainer2d.py (Rm, Lm; r_clip=0.05); ERRATA.md E-25 (worst |r_neg|/R = 0.054; r_clip = 0.05)
% TODO: r_clip = 0.05 is the Config default; not logged.

\paragraph{Loss function.}
Every physics-informed family minimises
\begin{equation}
  \mathcal{L} = \mathcal{L}_{\mathrm{data}}
              + \lambda_{\mathrm{bc}}\,\mathcal{L}_{\mathrm{bc}}
              + \lambda_{\mathrm{phys}}\,g\,\mathcal{L}_{\mathrm{phys}},
  \label{eq:loss}
\end{equation}
% src: trainer.py fit() (loss = data + lambda_bc·bc + lambda_physics·phys_gain·phys); trainer2d.py _loss_with_physics
with $\mathcal{L}_{\mathrm{data}}$ the mean squared error in standardised output
space. The boundary term enforces the traction-free condition at the free
surface,
\begin{equation}
  \mathcal{L}_{\mathrm{bc}} = \Bigl\langle
     \bigl(\hat\sigma_{rr}(r{=}1)/s_{rr}\bigr)^2
   + \bigl(\hat\tau_{rz}(r{=}1)/s_{rz}\bigr)^2 \Bigr\rangle_{\mathrm{batch}},
  \label{eq:lbc}
\end{equation}
% src: trainer.py _bc_loss (r1 = 1, p / ys); ERRATA.md E-01 (free surface at r = 1, condition and audit at the same point)
where $s_c$ is the training standard deviation of component $c$; the condition is
imposed and later audited at the same point, $r=1$ (at every $z$ in the extended
track). The physics term is a $\log(1+x^2)$ transform of the dimensionless
residual,
\begin{equation}
  \mathcal{L}_{\mathrm{phys}} = \Bigl\langle \log\!\bigl(1 + (R/s_R)^2\bigr) \Bigr\rangle,
  \label{eq:lphys}
\end{equation}
% src: trainer.py _physics_strong (torch.log1p((res/res_scale)**2).mean()); trainer2d.py (same form)
evaluated at the 20 data radii (collocation points coincide with data points),
% FIXED: source — trainer.py _physics_strong evaluates the residual at r_b (the N_R = 20 data radii); protocol.py (n_collocation=20) is the legacy six-model protocol and is not imported by the v2 launcher.
with scales $s_{R_1}=(s_{rr}+s_{\theta\theta})/R_{\mathrm{char}}$ for the radial and
$s_{R_z}=s_{zz}/L_{\mathrm{char}}+s_{rz}/R_{\mathrm{char}}$ for the axial equation,
$R_{\mathrm{char}}$ and $L_{\mathrm{char}}$ being the mean surface radius and window
length of the training set.
% src: trainer.py (res_scale = (sd_rr + sd_tt)/R_char, R_char = mean r_phys[:, -1] over training sets); trainer2d.py (res_r_scale, res_z_scale = sd_zz/L_char + sd_rz/R_char, L_char)
Without this scaling the physics term falls several orders of magnitude below
$\mathcal{L}_{\mathrm{data}}$ and is inert whatever $\lambda$; scaling the axial
residual with $R_{\mathrm{char}}$ instead of $L_{\mathrm{char}}$ has the same effect.
% src: trainer.py comment ("на пять порядков ниже data-loss"); PLAYBOOK.md §7.0.1 (same); trainer2d.py comment (axial residual scaled by R_char made the term numerically inert); PLAYBOOK.md §7.0.1a (L_char computed but unused — fixed)

The weights are fixed a priori and shared by all physics-informed families,
$\lambda_{\mathrm{phys}}=0.3$, $\lambda_{\mathrm{bc}}=0.5$,
% src: runs_v2.csv (lambda_physics=0.3 and lambda_bc=0.5 in all 234 pinn/vpinn/pinn2d_* rows; 0.0 in all 107 mlp/mlp2d rows)
but $\lambda_{\mathrm{phys}}$ is read as a share of the data gradient: before
training, $g$ is set once on the first mini-batch so that the gradient norms of
the physics and data terms with respect to the network parameters coincide,
\begin{equation}
  g = \lVert\nabla_{\theta}\mathcal{L}_{\mathrm{data}}\rVert_2 \big/
      \lVert\nabla_{\theta}\mathcal{L}_{\mathrm{phys}}\rVert_2 .
  \label{eq:calib}
\end{equation}
% src: trainer.py _calibrate, _grad_norm (idx0 = first batch_sets training cases, once before the epoch loop); trainer2d.py _calibrate
This balancing \cite{TODO-pinn-gradient-balancing} is what makes strong and weak
forms comparable: the weak residual projects the same residual onto oscillating
test functions and is damped roughly two-hundred-fold, so that at equal nominal
$\lambda$ the weak term contributed 0.36\,\% of the data-gradient norm against
37\,\% for the strong term, and ``strong versus weak'' degenerated into ``physics
on versus off''.
% src: PLAYBOOK.md §7.0.1a item 2 (0.36 % vs 37 %; damped ≈200-fold); trainer.py _calibrate docstring (≈200)
The measured multipliers in the main stage are $g\approx0.165$ (0.13--0.22) for
the strong and $g\approx4.97$ (2.8--8.0) for the weak form; $g$ is logged with
every PyTorch run.
% FIXED: scoped to the main stage; "every run" → "every PyTorch run" (the JAX twin does not calibrate, see Backends).
% src: RESULTS.md §1.4 (×0.165, ×4.97); runs_v2.csv phys_gain, stage=main: pinn 0.12710–0.22025 (mean 0.16499), vpinn 2.75588–8.02296 (mean 4.96797). Over all 1D PyTorch stages: pinn 0.127–0.256, vpinn 2.76–10.58 (corrupt stage widest). PLAYBOOK §7.0.1a quotes an earlier measurement (0.162 / 4.87) — not used.

\paragraph{Weak form (VPINN).}
Multiplying \eqref{eq:eq_red} by $r\,v_k(r)$ with test functions
$v_k=\sin(k\pi r)$, $k=1,\dots,8$, $v_k(0)=v_k(1)=0$, and integrating by parts over
$[0,1]$ gives \cite{TODO-vpinn}
\begin{equation}
  W_k = -\int_0^1 \hat\sigma_{rr}\,\bigl(v_k + r\,v_k'\bigr)\,dr
        + \int_0^1 \bigl(\hat\sigma_{rr}-\hat\sigma_{\theta\theta}\bigr)\,v_k\,dr ,
  \label{eq:weak}
\end{equation}
% src: trainer.py _physics_weak docstring and code (term1, term2; second integral carries the factor r·R/r_safe, equal to 1 away from the clipped axis region); Config.n_test_funcs=8 (default; not logged)
evaluated by 32-point Gauss--Legendre quadrature and inserted into
\eqref{eq:lphys} with scale $s_W=s_{rr}+s_{\theta\theta}$.
% src: trainer.py (Config.n_quadrature=32 default, _gauss_legendre mapped to [0,1]; weak_scale = sd_rr + sd_tt)
% TODO: n_test_funcs=8 and n_quadrature=32 are Config defaults; not logged.
No derivative of the network is computed: this, not another way of computing the
same derivative, is the substantive difference between the two families.
% src: trainer.py _physics_weak docstring; PLAYBOOK.md §7.0.1 table

\begin{table}[t]
% REV К12: "budget" в подписи раскрыт (эпохи, ранняя остановка, батчи)
% REV К12 (правка ревизии): "All share ... the seeds and the splits" неверно для шести семейств сразу — расширенный трек (mlp2d, pinn2d_*) прогнан на сидах 0–2 и на двух сплитах (runs_v2.csv, stage twod: random, extrap:alpha_max) против сидов 0–4 и тринадцати сплитов основного трека. Область действия ограничена треком.
\caption{Model families. Within each track all families share the network, the
optimiser, the training conditions (epoch limit, early-stopping rule,
mini-batch size), the seeds and the splits; only the loss composition differs. $R_1$: reduced radial residual
\eqref{eq:eq_red}; $R_r,R_z$: full residuals
\eqref{eq:eq_full_r}--\eqref{eq:eq_full_z}; $W_k$: weak residual \eqref{eq:weak}.}
\label{tab:families}
\centering\small
\begin{tabular}{@{}llllc@{}}
\toprule
Family & Input & Physics term & BC term & Network derivative \\
\midrule
MLP                     & $(Q,k,\alpha,\mu,v,r)$   & --                                        & --  & none \\
PINN (strong)           & $(Q,k,\alpha,\mu,v,r)$   & $R_1$, pointwise                          & yes & $\partial/\partial r$ (autodiff) \\
VPINN (weak)            & $(Q,k,\alpha,\mu,v,r)$   & $W_k$, $k\le8$, 32-pt quadrature          & yes & none \\
\texttt{mlp2d}          & $(Q,k,\alpha,\mu,v,z,r)$ & --                                        & --  & none \\
\texttt{pinn2d\_reduced}& $(Q,k,\alpha,\mu,v,z,r)$ & $R_r$ without $\partial\tau_{rz}/\partial z$ & yes & $\partial/\partial r$ \\
\texttt{pinn2d\_full}   & $(Q,k,\alpha,\mu,v,z,r)$ & $R_r$ and $R_z$, full                     & yes & $\partial/\partial r,\ \partial/\partial z$ \\
\bottomrule
\end{tabular}
% src: trainer.py FAMILIES and fit() (BC + physics added only for pinn/vpinn), trainer2d.py FAMILIES_2D (use_phys = family != mlp2d); runs_v2.csv lambda_bc = 0.5 for physics families, 0.0 for mlp/mlp2d
\end{table}

\paragraph{Training protocol.}
% REV К16: добавлена вводная конструкция перед протоколом ("Как отмечено выше, для корректного сравнения моделей необходимо обеспечить равные условия обучения, для чего использован следующий протокол")
As noted above, a meaningful comparison between the model families requires
that they be brought to the training stage under equal conditions; the protocol
described below was designed to that end.
The protocol is enforced by construction rather than by configuration: every
family is trained by the same trainer class, in which the family selects only the
composition of the loss, and one launcher writes every run with its
hyperparameters to a single log.
% FIXED: the draft cited protocol.py ("held in one module imported by every launcher; any deviation raises an assertion"). protocol.py is NOT imported by run_matrix.py, trainer.py or trainer2d.py (grep over publication/code and publication/experiments finds no import); it describes the legacy six-model/Optuna protocol and its constants contradict the v2 runs (PROTOCOL.max_epochs=400 vs 800 reached; PROTOCOL.patience=30 vs 40 for 1D); assert_conforms checks only optimizer, lr_schedule, n_trials, n_inner_splits, split_seed. Do not cite protocol.py for the v2 protocol.
% src: trainer.py (module docstring: one net, one optimiser, no scheduler, same epochs/batches/early stopping/seeds; Trainer.fit); trainer2d.py; run_matrix.py (one csv, FIELDS incl. lambda_*, learning_rate, n_units, n_layers, phys_gain); RESULTS.md header
It is identical for all families: AdamW \cite{TODO-adamw}, learning rate
$3\times10^{-3}$ ($2\times10^{-3}$ in the extended track), no schedule, weight
decay $10^{-4}$, gradient-norm clipping at 1.0, mini-batches of 64 cases (32),
at most 800 epochs (300), early stopping on the macro-RMSE of an internal
early-stopping subset of 15\,\% of the training cases with patience 40 (30), and % EDIT: "validation subset" → "early-stopping subset" (critic; keeps "validation" for comparison against experiment, as in §2.2)
restoration of the best state.
% src: runs_v2.csv (learning_rate 0.003 in all 323 1D rows, 0.002 in all 18 twod rows; epochs reaches 800 in main/curves/backend and exactly 300 in twod → max_epochs 800 / 300 were passed via --max-epochs); trainer.py fit (torch.optim.AdamW, no scheduler, clip_grad_norm_, best_state restored); jax_twin.py (optax.adamw + clip_by_global_norm); run_matrix.py one_run / one_run_2d: patience = job.get('patience', 40) / 30 — jobs never set it and an hp key 'patience' would raise a duplicate-kwarg TypeError (no run has an error), so 40 / 30 is confirmed.
% TODO (not logged, code defaults only; the --hp JSON is absent from the repository): weight_decay=1e-4, grad_clip=1.0 (Config, Config2D), batch_sets=64 (Config) / 32 (Config2D), val_frac=0.15. Confirm with the run author that --hp overrode only n_layers, n_units, learning_rate and lambda_bc (the four fields whose logged values differ from the Config defaults 4 / 128 / 1e-3 / 0.5).
% REV К14: утверждение о фиксированной архитектуре и гиперпараметрах перенесено в конец подраздела (см. абзац "Fixed architecture and hyperparameters" ниже)
Every configuration is trained with five seeds (0--4) on the thirteen splits of
Table~\ref{tab:splits} and in the backend comparison, and with three seeds (0--2)
in the learning-curve, label-corruption and extended-track stages;
% src: runs_v2.csv (seed ∈ {0..4} for stage main and backend; {0,1,2} for curves, corrupt, twod; 13 distinct split values)
the seed governs initialisation, batch order and the internal early-stopping subset, % EDIT: "validation subset" → "early-stopping subset" (critic)
while the test cases of a split are the same for every seed (split seed 42).
% src: trainer.py (torch.manual_seed(cfg.seed); rng = default_rng(cfg.seed) for the val subset; torch.Generator().manual_seed(cfg.seed) for batch order); splits.py build_split_suite(seed=42) called by run_matrix.py with the default; protocol.py split_seed=42 (legacy, consistent)
Learning curves use nested group-aware subsamples of 10/25/50/100\,\% of the
training pool, identical for all families;
% src: learning_curves.py DEFAULT_FRACTIONS=(0.10,0.25,0.50,1.00), nested_subsamples; run_matrix.py stage curves (seed=42); runs_v2.csv frac ∈ {0.1,0.25,0.5,1.0}
the subsamples are nested but not stratified by $v$ (the launcher uses
\texttt{nested\_subsamples}, not the stratified variant), so the sparsest level
$v = 5$~m/min (122 cases) can be under-represented at 10\,\%;
% EDIT: E-12 caveat added (critic). src: run_matrix.py l.188/202 (from learning_curves import nested_subsamples); learning_curves.py (stratified_nested_subsamples exists, unused by the launcher); ERRATA.md E-12 (122 cases at v = 5)
% TODO: report the count of v = 5 cases in the 10 % and 25 % subsamples (not logged in runs_v2.csv)
label corruption adds $+250$\,MPa to $\sigma_{\theta\theta}$ of 0/1/5/10\,\% of
the training cases (0\,\% being the clean reference; nested lists shared by all
families, hold-out clean), the morphology of a solver artefact found in the
source rather than Gaussian noise, which averages out over a profile.
% FIXED: added the 0 % reference (corruption.DEFAULT_RATES = (0.0, 0.01, 0.05, 0.10); runs_v2.csv corrupt_rate ∈ {0,0.01,0.05,0.1}; RESULTS.md §4)
% src: corruption.py (DEFAULT_OFFSET_MPA=250.0, DEFAULT_COMPONENT=1 → σ_θθ, make_corruption_plans nested, seed=42; docstring: artefact #439, whole set corrupted, hold-out clean); run_matrix.py (only y_tr corrupted)

% EDIT: the "Evaluation splits" table (a second \label{tab:splits}) deleted — the full 13-split table with n_train/n_test is Table~\ref{tab:splits} in Section~\ref{sec:fem-dataset}; its caption sentence on the 13.2 % friction-label duplicates was moved to §2.2 "Factorial design" (critic).

\paragraph{Metrics and physical audit.}
The primary metric is the macro-RMSE (mean over the four components of the
component RMSE, MPa). $R^2$ is reported with a single denominator, the standard
deviation of each component over the whole dataset, since a denominator taken
from the held-out slice is not comparable between regions.
% FIXED: the draft stated that per-component NRMSE also uses the global denominator. In RESULTS.md §1.3б the per-component NRMSE table is normalised by the held-out std (make_results.py: "NRMSE = RMSE / std компоненты на held-out", columns nrmse_*), whereas R² in §1.0б is global (add_global_metrics). Global per-component columns (nrmse_*_global) exist in runs_v2.csv but are not tabulated.
% TODO: decide one denominator for Table 2 (recommended: global, columns nrmse_*_global / r2_*_global) and say so in the caption.
% src: trainer.py metrics() docstring (local vs global; macro_rmse = mean of component RMSE); make_results.py add_global_metrics; RESULTS.md §1.0б, §1.3б
Physical admissibility is audited on the predictions with the operators of the
loss (median $|R_1|$ over the profile by finite differences; $|\hat\sigma_{rr}|$,
$|\hat\tau_{rz}|$ at $r=1$), always beside the same quantity computed on the FEM
data of the same test cases, since the source itself satisfies the traction-free
condition only approximately.
% src: trainer.py equilibrium_audit (np.gradient, median |res| over interior points, then median over cases; r_clip 0.05), bc_audit; run_matrix.py (bc_audit at predict_at(r=1.0)); RESULTS.md §1.4 (FEM baseline rows: |σ_rr| 3.52, |τ_rz| 0.24 MPa overall)

\paragraph{Backends.}
The reference implementation is PyTorch (\texttt{torch.autograd}); an independent
twin in JAX (\texttt{jax.grad}, \texttt{optax.adamw}) shares architecture, data,
split, the epoch limit and early-stopping rule, and the nominal hyperparameters,
% REV К12: "training budget" заменён на явное перечисление (предел числа эпох и правило ранней остановки)
but draws initial weights and batch order from its own random stream.
% src: jax_twin.py module docstring; JaxTrainer.fit (optax.chain(clip_by_global_norm, adamw); batch order default_rng(seed + 10000); same val_frac split rule)
The twin implements the strong-form loss with the nominal weight
$\lambda_{\mathrm{phys}}$ applied directly: the gradient-norm calibration of
\eqref{eq:calib} is not implemented in it, so the effective physics weight of the
two implementations differs.
% FIXED: the draft claimed the twin "shares ... protocol ... hyperparameters" without qualification. jax_twin.py JaxTrainer.fit has no _calibrate / phys_gain (loss = data + lam_bc·bc + lam_p·log1p(...) with lam_p = cfg.lambda_physics = 0.3); runs_v2.csv logs phys_gain = 1.0 for all 10 backend/jax rows against 0.132–0.215 for the 10 backend/torch rows. Effective weight: 0.3 (JAX) vs 0.3 × 0.165 ≈ 0.05 (PyTorch).
% TODO: either port _calibrate to jax_twin.py and rerun the 10 JAX runs, or state this difference explicitly in §3.2(2) — it is a confound of the "implementation" comparison.
With identical weights copied into both networks, $\partial\hat\sigma_{rr}/\partial r$
at 64 random inputs agrees to $8.7\times10^{-15}$ absolute and $9.8\times10^{-16}$
relative (forward pass $1.0\times10^{-15}$), i.e.\ to double-precision round-off,
so no accuracy difference between backends can originate in automatic
differentiation.
% src: PLAYBOOK.md §7.0.2 (forward 1.0e-15; gradient 8.7e-15, relative 9.8e-16); ERRATA.md E-19 (8.7e-15, 9.8e-16); jax_twin.py check_gradient_agreement (n_points=64, X ~ N(0,1), gradient w.r.t. column 5 = r)
Each backend is then trained with five seeds on \texttt{random} and
\texttt{extrap:alpha\_max} and compared by the paired test below, with wall-clock
training time measured on CPU in double precision.
% FIXED: "one CPU thread" → "CPU". run_matrix.py torch.set_num_threads(1) constrains PyTorch only; jax_twin.py sets no XLA thread flags.
% TODO: establish the thread count actually used by the JAX runs before writing "one thread"; RESULTS_v1_vs_v2.md §7 states only "CPU, float64" (313 ± 102 s vs 516 ± 94 s, ×1.65 — consistent with runs_v2.csv seconds, stage backend).
% src: runs_v2.csv (stage backend: backend ∈ {torch, jax}, seeds 0–4, splits random and extrap:alpha_max, 20 rows)
% REV К17: отсылка к предыдущей версии текста снята; утверждение сформулировано самостоятельно
% REV К17 (правка ревизии): "DeepXDE ... runs on the same PyTorch backend" как свойство библиотеки источником не подтверждается (DeepXDE поддерживает и другие бэкенды). ERRATA E-19 фиксирует конкретное: models/pinn_model_dde.py задаёт DDE_BACKEND=pytorch, и dde.grad.jacobian вызывает torch.autograd.grad. Формулировка возвращена к тому, что есть в источнике.
The DeepXDE implementation in the same code base \cite{TODO-deepxde} is
configured with the PyTorch backend and obtains its derivatives through
\texttt{torch.autograd.grad}; it is therefore not an independent
differentiation engine and is not reported as a third backend.
% src: ERRATA.md E-19 (DDE_BACKEND=pytorch, dde.grad.jacobian → torch.autograd.grad); jax_twin.py docstring

\paragraph{Statistical treatment.}
The unit of replication is the training run (seed), not the test case. Within a
split, two families are compared by a paired $t$-test over seeds ($\alpha=0.05$;
no verdict with fewer than three seeds), with $|\Delta|$ reported beside the pooled
seed standard deviation as effect size.
% src: stats.py significance_gate (scipy ttest_rel, alpha=0.05, MIN_SEEDS_FOR_TEST=3, pooled = sqrt(0.5(s_i² + s_j²)), practical flag |Δ| ≥ k_sigma·pooled)
Ranking of the three families within a split uses Friedman with Nemenyi post hoc,
blocks being seeds \cite{TODO-demsar-2006}; indistinguishability is not
transitive, so distinguishable pairs are listed rather than ``clusters''.
% src: stats.py friedman_test (n_blocks ≥ 3), nemenyi_cd, NemenyiResult.cliques / to_text ("различимые пары"); ERRATA.md E-15
Splits are never pooled: ranks in \texttt{extrap:alpha\_max} and in
\texttt{random} describe different tasks.
% src: stats.py pool_guard
The excess error of an extrapolation split over its matched control is tested by
Welch's unpaired $t$-test; trends across data fraction and corruption rate by
least squares of the seed-wise PINN$-$MLP gap on $\log$(fraction) and on the rate.
% src: stats.py delta_vs_matched (ttest_ind, equal_var=False); paper_tables.py (scipy.stats.linregress of the seed-level PINN−MLP gap on corrupt_rate); RESULTS_v1_vs_v2.md §5 (regression on log(доля)) and §6 (slope −14.20 MPa per unit fraction, p = 0.00036, R² = 0.735)
% REV К17: отсылка к предыдущей версии текста заменена самостоятельным утверждением об отдельном прогоне тестов; числа (304 hold-out случая) не изменены
A separate run of the tests, in which each model was trained once and
Friedman--Nemenyi was applied across the 304 hold-out cases, shows what such a
design measures: the sampling variance of the hold-out, that is, on how many
cases one model beats another, and not whether an advantage reproduces on
retraining, which is the question at issue here.
% src: ERRATA.md E-08 (one run per model; Friedman–Nemenyi over 304 hold-out sets); stats.py module docstring (304)

\paragraph{Fixed architecture and hyperparameters.}
% REV К14: утверждение перенесено сюда из абзаца "Training protocol" (руководитель просил поместить его в более позднюю часть текста)
% TODO К14: согласовать с руководителем, описывать ли подбор Optuna из предшествующей работы и где. В коде настоящей публикации Optuna не используется (publication/code/trainer.py, protocol.py: protocol.py с полями n_trials/Optuna не импортируется ни run_matrix.py, ни trainer.py, ни trainer2d.py), поэтому приписывать текущим прогонам оптимизацию Optuna нельзя.
% TODO: state how the 3×96 / 3e-3 configuration was chosen — no "tune" stage appears in runs_v2.csv (stages: main 195, curves 72, corrupt 36, backend 20, twod 18); the Config defaults are 4×128 / 1e-3, so the choice was made outside the logged matrix.
One architecture and one hyperparameter set were fixed before any held-out
region was evaluated and were not tuned per family or per split, so no search
could leak a held-out region.
% src: RESULTS.md header ("Гиперпараметры зафиксированы априори и на held-out регионах не подбирались"); PLAYBOOK.md §8 item 1 ("Optuna на held-out регионе не запускается вовсе").
% CHECKS: (a) no data-deficit claim is made (learning curves described as protocol only; RESULTS_v1_vs_v2.md §5: thesis NOT confirmed); (b) Johnson–Cook constants are not mentioned; (c) the only "validation" is the internal early-stopping subset (ML sense) — no FEM "validation" is claimed (E-09: none exists); (d) the 2.2/2.3 assumption comment stands at the top.